\section{Jobans Abschnitt}
\abschnittsautor{J. Rahal}


\subsection{Systemarchitektur}
Webbasierte Systeme beruhen auf dem Client-Server Modell. Ein Client (Webrowser)
stellt Anfragen, ein Server nimmt diese entgegen, verarbeitet sie und sendet eine Antowrt zurück. 
Die Aufgaben werden hier auf drei Schichten verteilt. Die Präsentationsschicht
ist zuständig für die Darstellung und die Interaktion mit dem Benutzer. Die Logikschicht verarbeitet die Daten
und die Datenhaltungsschicht speichert sie dauerhaft. 

Der Vorteil daran ist, dass jede Schicht unabhängig voneinander weiterentwickelt oder gewechselt werden kann,
solange die Schnittstellen unverändert bleiben. Für VerTex ist das
wesentlich, weil Bildverarbeitung, Backend und Weboberfläche von
unterschiedlichen Personen parallel entwickelt werden.


\subsection{Kommunikation zwischen Frontend und Backend}

